\subsection{Associated coverings}

Let B be a topological space, let G be a discrete topological group, and let E be a principal covering of B with group G. Let F be a G-set; if F is endowed with the discrete topology, the group G acts continuously on F. The B-space \(E \times^{G} F\) is then a covering. It is called the covering of B with fiber type F associated with the principal covering E.

DEFINITION 5. --- Let B be a topological space, let G be a discrete topological group, and let E be a principal covering of B with group G. A covering X of B is said to be associable to the principal covering E if there exists a G-set F and a B-isomorphism from E × \(^{G}\) F onto X.

Let B be a topological space, let G be a discrete topological group, and let E be a principal covering of B with group G.

For every covering X of B, the group G acts on the left on \(\mathcal{C}_{\mathrm{B}}(\mathrm{E};\mathrm{X})\) by the law defined by \((g\cdot h)(x)=h(x\cdot g)\) for \(x\in E\) , \(g\in G\) , \(h\in\mathcal{C}_{\mathrm{B}}(\mathrm{E};\mathrm{X})\) .

For every G-set F equipped with the discrete topology, define a map \(\alpha_{\mathrm{F}}\colon \mathrm{F}\to \mathcal{C}_{\mathrm{B}}(\mathrm{E};\mathrm{E}\times^{\mathrm{G}}\mathrm{F})\) by:

\[
\alpha_ {\mathrm{F}} (y) (x) = \pi (x, y) \quad \text {for } y \in \mathrm{F} \text { and } x \in \mathrm{E},\tag{2}
\]

where \(\pi\colon\mathrm{E}\times\mathrm{F}\to\mathrm{E}\times^{\mathrm{G}}\mathrm{F}\) is the canonical surjection. The map \(\alpha_{\mathrm{F}}\) is compatible with the operations of G on F and on \(\mathcal{C}_{\mathrm{B}}(\mathrm{E};\mathrm{E}\times^{\mathrm{G}}\mathrm{F})\) .

For every covering X of B, there exists a unique map \(\beta_{X}\) from \(E \times^{G} \mathcal{C}_{B}(E; X)\) into X such that:

\[
\beta_ {\mathrm{X}} (\pi (x, h)) = h (x) \quad \text {for } h \in \mathcal {C} _ {\mathrm{B}} (\mathrm{E}; \mathrm{X}) \text { and } x \in \mathrm{E}.\tag{3}
\]

Indeed, we have \((g^{-1}\cdot h)(x\cdot g)=h(x)\) for \(x\in E\) , \(g\in G\) and \(h\in\mathcal{C}_{\mathrm{B}}(\mathrm{E};\mathrm{X})\) , by definition of the action of G on \(\mathcal{C}_{\mathrm{B}}(\mathrm{E};\mathrm{X})\). When one equips the set \(\mathcal{C}_{\mathrm{B}}(\mathrm{E};\mathrm{X})\) with the discrete topology, the map \(\beta_{X}\) is a B-morphism of coverings.

When \(\mathrm{X} = \mathrm{E}\times^{\mathrm{G}}\mathrm{F}\) , we have

\[
\beta_ {\mathrm{X}} \left(\pi \left(x, \alpha_ {\mathrm{F}} (y)\right)\right) = \pi (x, y) \quad \text {   for   } x \in \mathrm{X} \text {   and   } y \in \mathrm{F}.\tag{4}
\]

PROPOSITION 8. --- Let B be a connected and nonempty topological space. Let G be a discrete group and let (E, p) be a principal covering of B with group G. Suppose that E is connected. With the notations above, we have:

\begin{enumerate}
\def\labelenumi{\alph{enumi})}
\item
  For every G-set F equipped with the discrete topology, the map \(\alpha_{F}\) is an isomorphism of the G-set F onto the G-set \(\mathcal{C}_{\mathrm{B}}(\mathrm{E};\mathrm{E}\times^{\mathrm{G}}\mathrm{F})\) .
\item
  Let X be a covering of B. The B-morphism \(\beta_{X}\) is an isomorphism from \(E \times^{G} \mathcal{C}_{B}(E; X)\) onto X if and only if the covering \((E \times_{B} X, pr_{1})\) of E is trivializable.
\item
  Let \(y\) and \(y'\) be points of F such that \(\alpha_{\mathrm{F}}(y) = \alpha_{\mathrm{F}}(y')\). The space E is not empty; let us choose a point \(x\) in it. We have \(\pi(x, y) = \pi(x, y')\) in \(\mathrm{E} \times^{\mathrm{G}} \mathrm{F}\). Consequently, there exists \(g \in \mathrm{G}\) such that \(x \cdot g = x\) and \(g^{-1} \cdot y = y'\). The first equality implies that \(g\) is the identity element \(e\) of G, hence \(y = y'\). Thus, the map \(\alpha_{\mathrm{F}}\) is injective. On the other hand, let \(h \in \mathcal{C}_{\mathrm{B}}(\mathrm{E}; \mathrm{E} \times^{\mathrm{G}} \mathrm{F})\) and let \(x\) be a point of E. Let \(x' \in \mathrm{E}\) and \(y' \in \mathrm{F}\) be such that \(h(x) = \pi(x', y')\). In particular, \(p(x) = p(x')\); there then exists an element \(g\) of G such that \(x' = x \cdot g\), and one also has \(h(x) = \pi(x, y)\), where we have set \(y = g \cdot y'\). The B-morphisms \(h\) and \(\alpha_{\mathrm{F}}(y)\) coincide at the point \(x\) of E; they are therefore equal since the space E is connected (I, p. 34, cor. 1 of prop. 11), and this proves that the map \(\alpha_{\mathrm{F}}\) is surjective.
\item
  By proposition 7, b) of I, p. 105, applied to F = \(\mathcal{C}_{\mathrm{B}}(\mathrm{E};\mathrm{X})\), the covering \(p^{*}(\mathrm{E}\times^{\mathrm{G}}\mathcal{C}_{\mathrm{B}}(\mathrm{E};\mathrm{X}))\) of E is isomorphic to the trivial covering \(\mathrm{E}\times \mathcal{C}_{\mathrm{B}}(\mathrm{E};\mathrm{X})\). If \(\beta_{\mathrm{X}}\) is an isomorphism, the covering \(p^{*}(\mathrm{X})\) of E is therefore trivializable. Conversely, suppose the covering \(p^{*}(\mathrm{X})\) is trivializable and let us show that the B-morphism \(\beta_{\mathrm{X}}\) is bijective; it will follow that \(\beta_{\mathrm{X}}\) is a B-isomorphism (I, p. 30, cor. 2 of prop. 6).
\end{enumerate}

The map \(\beta_{\mathrm{X}}\) is obtained by passing to the quotient from the map \(\gamma\colon\mathrm{E}\times\mathcal{C}_{\mathrm{B}}(\mathrm{E};\mathrm{X})\to\mathrm{X}\) defined by \(\gamma(x,h)=h(x)\) . Let us show that the map \(\gamma\) is surjective. Let \(x\) be a point of X. The projection of the B-space E is surjective, since it is a principal covering; there therefore exists a point \(y\) of E such that \((y,x)\in\mathrm{E}\times_{\mathrm{B}}\mathrm{X}\) . There then exists a continuous section \(s\) of the trivializable covering \(\mathrm{pr}_{1}:\mathrm{E}\times_{\mathrm{B}}\mathrm{X}\to\mathrm{E}\) such that \(s(y)=(y,x)\) . The map \(h=\mathrm{pr}_{2}\circ s:\mathrm{E}\to\mathrm{X}\) is a B-morphism and one has \(h(y)=x\) , whence the surjectivity of the map \(\gamma\) and, consequently, that of the map \(\beta_{\mathrm{X}}\) .

Let us finally show that \(\beta_{\mathrm{X}}\) is injective. Let \((x,h)\) and \((x^{\prime},h^{\prime})\) be elements of \(\mathrm{E}\times \mathcal{C}_{\mathrm{B}}(\mathrm{E};\mathrm{X})\) such that \(h(x) = h'(x')\) . Note that \(x\) and \(x'\) have the same projection in B; there therefore exists an element \(g\) of G such that \(x' = x\cdot g\) . One then has \(h(x) = h'(x\cdot g) = (g\cdot h')(x)\) . Since the space E is connected, one has \(h = g\cdot h'\) (I, p. 34, cor. 1 of prop. 11). Thus, \((x,h)\) and \((x',h')\) have the same class in \(\mathrm{E}\times^{\mathrm{G}}\mathcal{C}_{\mathrm{B}}(\mathrm{E};\mathrm{X})\) , which shows that the map \(\beta_{\mathrm{X}}\) is injective, which completes the proof.

COROLLARY 1. --- Let (E,p) be a principal covering of B with group G; suppose that E is connected and nonempty. A covering X of B is associable to E if and only if the covering \(p^{*}(X)\) is trivializable.

If the covering \(p^{*}(\mathrm{X})\) is trivializable, it follows from Proposition 8, b), that the covering X is associable to E. In this case, the map \(\beta_{X}\) identifies the covering X with the fiber-type covering \(\mathcal{C}_{\mathrm{B}}(\mathrm{E};\mathrm{X})\) associated with E. Conversely, let F be a G-set equipped with the discrete topology, and suppose that we have \(X = E \times^{G} F\) . Then, \(\alpha_{F}\) is an isomorphism (loc. cit., a)) and formula (4) implies that \(\beta_{X}\) is an isomorphism. Consequently, the covering \(p^{*}(\mathrm{X})\) is trivializable (prop. 8, b)), whence the corollary.

COROLLARY 2. --- Let B be a connected and locally connected topological space, let (E,p) be a principal covering of B with group G. Let E₀ be a connected component of E and let G₀ be the subgroup of G stabilizing E₀. The B-space (E₀, p \textbar{} E₀) is a principal covering with group G₀ (I, p. 103, prop. 6).

Every covering X of B that is associable to the principal covering E is associable to the principal covering \(E_{0}\). In particular, the covering E is associated with the principal covering \(E_{0}\).

Indeed, if the covering X is associable to E, the covering \(p^{*}(\mathrm{X})\) is trivializable and the covering \(p_{0}^{*}(\mathrm{X})\) induced by \(p^{*}(\mathrm{X})\) over \(E_{0}\) is therefore trivializable.

More precisely, note that the map \((x,g)\mapsto x\cdot g\) from \(E_{0}\times G\) into E induces, by passing to the quotient, a B-isomorphism of principal coverings from \(E_{0}\times^{G_{0}}G\) onto E.

PROPOSITION 9. --- Let B be a topological space, let E be a principal covering of B with group G, connected and nonempty. Let F be a nonempty G-set equipped with the discrete topology. For the space E × \(^{G}\) F to be connected, it is necessary and sufficient that G acts transitively on F.

Let U be an open and closed subset of \(E \times^{G} F\) . If \(\pi: E \times F \to E \times^{G} F\) denotes the canonical surjection, \(\pi^{-1}(U)\) is an open and closed subset of \(E \times F\) which is stable under G. Since E is connected, there exists a subset \(F' \subset F\), stable under G, such that \(\pi^{-1}(U) = E \times F'\), whence \(U = \pi(E \times F')\). Let \(F'\) and \(F''\) be subsets of F stable under G such that \(\pi(E \times F') = \pi(E \times F'')\). Since E is not empty, we have \(F' = F''\). The map \(F' \mapsto \pi(E \times F')\) is a bijection from the set of G-stable subsets of F onto the set of clopen subsets of \(E \times^{G} F\). The proposition follows.

PROPOSITION 10. --- Let B be a topological space, let (E,p) be a principal covering of B with group G, and let X be a covering of B. Suppose that E and X are connected and nonempty. The following properties are equivalent:

\begin{enumerate}
\def\labelenumi{(\roman{enumi})}
\item
  The covering X is associated to the principal covering E;
\item
  There exists a subgroup H of G such that X is B-isomorphic to E/H;
\item
  There exists a surjective B-morphism \(h: \mathrm{E} \to \mathrm{X}\) ;
\item
  For every point \((y, x)\) of \(E \times_{B} X\) there exists a B-morphism \(r: E \to X\) such that \(r(y) = x\) .
\end{enumerate}

Suppose these conditions are satisfied, and let H be a subgroup of G such that X is B-isomorphic to E/H. The covering X is Galois if and only if the subgroup H is normal in G.

(i)⇒(ii): Let F be a discrete G-set such that the covering \(E \times^{G} F\) is B-isomorphic to X. The set F is nonempty and the space \(E \times^{G} F\) is connected, hence the group G acts transitively on F (prop. 9). Then, the space \(E \times^{G} F\) is B-isomorphic to E/H, where H is the subgroup of G fixing a point of F (I, p. 106, example 4).

(ii)⇒(iii): Indeed, the canonical surjection from E onto E/H is a surjective B-morphism.

(iii)⇒(iv): Let \((y, x)\) be a point of \(E \times_{B} X\). Since the map h is surjective, there exists a point \(y'\) of E such that \(h(y') = x\). The points y and \(y'\) have the same projection in B. Since the covering E is principal with group G, there exists \(g \in G\) such that \(y \cdot g = y'\). The map \(r: E \to X\) defined by \(z \mapsto h(z \cdot g)\) is a B-morphism and one has \(r(y) = x\).

(iv)⇒(i): Since X is nonempty and the projection of E onto B is surjective, the space \(E \times_{B} X\) is not empty. There therefore exists a B-morphism r from E to X. The map ( \(Id_{E}, r\) ): \(E \to E \times_{B} X\) is a section of the covering \(p^{*}(X)\) of E. Since the space E is connected, it follows from Corollary 2 of Prop. 1 in I, p. 69 that the covering \(p^{*}(X)\) is trivializable, which proves that the covering X is associable to the Galois covering p (Proposition 8).

Suppose now that these conditions are satisfied. Let X denote the B-space E/H, q its projection, and \(h: E \rightarrow E/H\) the canonical map.

If the group H is normal in G, then X is a principal covering with group G/H (I, p. 106, example 4). Since X and B are connected and nonempty, X is a Galois covering of B (I, p. 102, th. 2). Conversely, suppose that E/H is a Galois covering of B and let K = Aut \(_{B}\) (E/H) \(^{\circ}\) . We define a map \(\alpha: \mathrm{E} \times \mathrm{G} \to \mathrm{K}\) by associating with \((t, g) \in \mathrm{E} \times \mathrm{G}\) the unique element \(k\) of K such that \(h(t \cdot g) = h(t) \cdot k\) . It is continuous because it is obtained by composing the continuous map \((t, g) \mapsto (h(t), h(t \cdot g))\) from E \(\times\) G into X \(\times_{\mathrm{B}}\) X with the canonical trivialization (I, p. 95, remark 1) X \(\times_{\mathrm{B}}\) X \(\to\) X \(\times\) K of \(q^{*}(\mathrm{X})\) . Since E is connected and K is a discrete group, the continuous map \(t \mapsto \alpha(t, g)\) is constant for every \(g \in \mathrm{G}\); we denote its value by \(\alpha(g)\). If \(t \in \mathrm{E}\), \(g \in \mathrm{G}\), and \(g' \in \mathrm{H}\), we then have

\[
h (t) = h \left(t \cdot g ^ {- 1} g\right) = h \left(t \cdot g ^ {- 1}\right) \cdot \alpha (g) = h \left(t \cdot g ^ {- 1} \cdot g ^ {\prime}\right) \cdot \alpha (g) = h \left(t \cdot \left(g ^ {- 1} g ^ {\prime} g\right)\right).
\]

It follows that \(g^{-1}g'g\) belongs to H, and therefore H is normal in G.

COROLLARY. --- Let E be a Galois covering of B and let X be a covering of B, connected and nonempty. Suppose that X is a finite covering or that the space B is locally connected. If there exists a B-morphism h: E → X, then the covering X is associated with the principal covering E.

In both cases, the B-morphism h is a covering (I, p. 77, cor. 3 of th. 1, and I, p. 78, prop. 5), and a nonempty covering of a connected space is surjective (I, p. 68).

THEOREM 3. --- Let B be a nonempty, connected, and locally connected topological space. Every covering of B can be associated with a Galois covering; every finite covering of B can be associated with a finite Galois covering.

Let X be a covering of B, and denote its projection by q. Since the space B is connected and nonempty, the covering X has a degree; denote this degree by F and endow the set F with the discrete topology. Since the space B is locally connected, the sheaf \(\mathcal{I} = \underline{\mathcal{I}som}_{\mathrm{B}}(\mathrm{B} \times \mathrm{F}; \mathrm{X})\) is locally constant, and at every point b of B, its stalk \(I_{b}\) is canonically isomorphic to the set \(\mathcal{B}(\mathrm{F}; \mathrm{X}_{b})\) of bijections from F onto the fiber \(X_{b}\) (I, p. 89, prop. 12). Let \(E = E_{\mathcal{I}}\) be the associated covering (I, p. 86, corollary).

Let G be the group of permutations of F. For every \(g \in G\), we define a morphism \(\gamma'(g)\) of the sheaf \(\mathcal{I}\) into itself as follows: for every open set U of B, the map \(\gamma'(g)_{\mathrm{U}}\) associates to a U-isomorphism \(\varphi: \mathrm{U} \times \mathrm{F} \to q^{-1}(\mathrm{U})\) the U-isomorphism defined by \((b, f) \mapsto \varphi(b, g(f))\) for \(b \in U\) and \(f \in F\). We have \(\gamma'(\mathrm{Id}_{\mathrm{F}}) = \mathrm{Id}_{\mathcal{I}}\) and \(\gamma'(g \circ g') = \gamma'(g') \circ \gamma'(g)\), for \(g, g' \in G\), so that for every \(g \in G\), \(\gamma'(g)\) is an automorphism of the sheaf \(\mathcal{I}\). The B-morphism \(\gamma(g): E \to E\) associated with \(\gamma'(g)\) is an automorphism (I, p. 50). If \(x = [\mathrm{U}, \varphi, b]\) is an element of E, where U is an open subset of B, b a point of U, and \(\varphi\) a U-isomorphism from U×F to \(q^{-1}(\mathrm{U})\), we have \(\gamma(g)(x) = [\mathrm{U}, \psi, b]\), where \(\psi\) is defined by \(\psi(a, f) = \varphi(a, g(f))\), for \(a \in U\) and \(f \in F\). If g and \(g'\) are elements of G, we have \(\gamma(g \circ g') = \gamma(g') \circ \gamma(g)\). We have \(\gamma(\mathrm{Id}_{\mathrm{F}}) = \mathrm{Id}_{\mathrm{E}}\). We thus define a continuous right action of G on E by setting \(x \cdot g = \gamma(g)(x)\), for \(x \in E\) and \(g \in G\). The group G acts simply transitively on every fiber of E, so that the covering E equipped with this action is a principal covering with group G (I, p. 97, prop. 4). Let h be the map from E × F to X defined by \(h([\mathrm{U}, \varphi, b], f) = \varphi(b, f)\) for every open subset U of B, every point b of U, and every U-isomorphism \(\varphi\) from U × F onto \(q^{-1}(\mathrm{U})\). By definition of the topology on E, it is continuous; it is a B-morphism. For every element g of G and every point \((x, f)\) of E × F, we have \(h(x, g(f)) = h(x \cdot g, f)\). The map h therefore defines, by passing to the quotient, a B-morphism \(h': E \times^{G} F \to X\). For every point \(b\) of B, the map \(h_b\colon \mathrm{E}_b\times \mathrm{F}\to \mathrm{X}_b\) is identified with the map \((\varphi, f)\mapsto \varphi(f)\) from \(\mathcal{B}(\mathrm{F};\mathrm{X}_b)\times \mathrm{F}\) into \(\mathrm{X}_b\), so that the map \(h_b^\prime\) is bijective. Consequently, \(h^\prime\) is a B-isomorphism from \(\mathrm{E}\times^{\mathrm{G}}\mathrm{F}\) onto X (I, p. 30, cor. 2 of prop. 6).

Since the space B is connected, locally connected, and nonempty, the covering X, associable to the principal covering E, is associable to a Galois covering (I, p. 110, corollary 2). If the covering X is finite, then so is the covering E, hence X is associable to a finite Galois covering (loc. cit.).

